% Table S8: Cohort and estimand separation
\begin{table}[htbp]
\centering
\caption{Questions addressed and interpretation boundaries. The Set-C docking cohort has 17 candidates and 136 wild-type/mutant systems after prioritization filters; 12 candidates have complete PfDHFR/PfCRT panels and define the primary RRS analysis, and 5 are PfCRT-only records reported separately. Parent-study MD systems do not overlap Set-C. The 16-system Set-C pilot passed geometric trajectory checks and is analyzed separately; its MD-RRS and MM-GBSA values are not used to assign docking-RRS classes.}
\label{tab:s6_cohort_estimands}
\scriptsize
\begin{tabularx}{\textwidth}{>{\raggedright\arraybackslash}p{0.16\textwidth} >{\raggedright\arraybackslash}p{0.17\textwidth} >{\raggedright\arraybackslash}p{0.13\textwidth} >{\raggedright\arraybackslash}X >{\raggedright\arraybackslash}X}
\toprule
Evidence layer & Cohort & Systems & Primary analyses & Interpretation boundary \\
\midrule
Set-C docking & PP-01--PP-17 & 136 (12 complete; 5 PfCRT-only) & WT/mutant docking; target-specific RRS, ACSI, PNS & Filtered prioritisation cohort; not measured binding or MD-RRS \\
Parent-study MD & Ligands 164, 201, 214, 438 & 4 WT complexes $\times$ \SI{10}{\nano\second} & RMSD, contacts, bound-state assessment, MM-GBSA where valid & Structural check; one interpretable MM-GBSA estimate (PfCRT--214) \\
Set-C MD pilot & PP-01/PP-02 & 16 systems (all met trajectory criteria) & Trajectory geometry, MD-RRS, MM-GBSA & Single-replicate structural follow-up; separate from docking-RRS classification \\
Full Set-C MD-RRS & PP-01--PP-17 & 136 planned systems & Full trajectory QC and MD-RRS & Not available because the complete cohort was not simulated \\
\bottomrule
\end{tabularx}
\end{table}
